\documentclass[11pt,letterpaper]{article}

\usepackage[margin=1in]{geometry}
\usepackage{amsmath}
\usepackage{graphicx}
\usepackage{booktabs}
\usepackage{hyperref}
\usepackage{caption}
\usepackage{subcaption}
\usepackage{float}
\usepackage{setspace}
\usepackage{parskip}
\usepackage{microtype}
\usepackage{xcolor}
\usepackage{enumitem}

\hypersetup{
    colorlinks=true,
    linkcolor=blue,
    citecolor=blue,
    urlcolor=blue
}

\title{%
  \textbf{Does Data Augmentation Benefit Scale With Architecture Depth?}\\[0.5em]
  \large A Systematic Study of CNN Depth, Augmentation Intensity,\\
  and Optimizer Choice on CIFAR-10\\[1em]
  \normalsize COGS 181 Final Project
}
\author{Jae Kim}
\date{March 2026}

\begin{document}

\maketitle
\thispagestyle{empty}

% ─────────────────────────────────────────────────────────────────────────────
\begin{abstract}
% ─────────────────────────────────────────────────────────────────────────────
Data augmentation is widely used to improve CNN generalization, but its interaction with model
capacity and optimizer choice is understudied. We investigate whether the benefit of increasingly
aggressive augmentation scales with architecture depth on CIFAR-10, and whether optimizer
selection further modulates this relationship. Two convolutional networks are compared: a
ShallowCNN (3 conv blocks, 620K parameters) and a DeepCNN (6 conv layers in a VGG-style
arrangement, 3.25M parameters), each trained under four augmentation levels (none, basic,
medium, aggressive) for 50 epochs. During the initial experiment we discovered that DeepCNN
fails to train at the standard learning rate of 0.1 with SGD, collapsing to random-chance
accuracy under any augmentation. We diagnose this as a learning rate instability driven by
augmentation-induced gradient variance in deep non-residual networks, and resolve it by
reducing the learning rate to 0.01 (SGD, 74.0\%). Switching to Adam ($\mathrm{lr}=10^{-3}$)
further boosts DeepCNN to \textbf{89.1\%} under basic augmentation---a +15.1 percentage point
gain over the tuned SGD baseline and the highest accuracy across all conditions. ShallowCNN is
more robust to aggressive augmentation ($-9$ pp) than DeepCNN ($-30$ pp), supporting
the hypothesis. Optimizer choice emerges as the single highest-impact hyperparameter in this study.
\end{abstract}

\setstretch{1.1}

% ─────────────────────────────────────────────────────────────────────────────
\section{Introduction}
% ─────────────────────────────────────────────────────────────────────────────

Data augmentation---artificially expanding the training distribution through random image
transformations---is one of the most reliable regularization techniques in deep learning.
By exposing the network to diverse variations of each training example, augmentation reduces
overfitting and can meaningfully improve test accuracy \cite{shorten2019survey}. However,
augmentation also adds variance to each gradient update, which may interact differently with
networks of different capacities, depths, and optimizers.

Our research question is: \textit{Does the benefit of aggressive data augmentation scale with
architecture depth on CIFAR-10?} The intuition motivating this question has two competing
sides. On one hand, deeper networks have greater capacity and thus more to gain from the
implicit regularization that augmentation provides. On the other hand, deeper non-residual
networks are harder to optimize and may be destabilized by the gradient noise that aggressive
augmentation introduces.

To study this, we design a controlled experiment on CIFAR-10 \cite{krizhevsky2009learning},
comparing a ShallowCNN (3 conv blocks) and a DeepCNN (6 conv layers, VGG-inspired
\cite{simonyanvery}) under four augmentation levels: none, basic (horizontal flip only), medium
(flip + crop + color jitter), and aggressive (medium + grayscale + random erasing). The
experiment produced two unexpected secondary findings. First, DeepCNN is highly sensitive to
learning rate at the standard value of 0.1 with SGD, failing to train under any augmentation
and achieving only 41.5\% even without augmentation. Second, after fixing the learning rate,
switching from SGD to Adam yields a dramatic +15 pp improvement, reaching 89.1\% on the best
configuration. Together, these findings show that hyperparameter interactions---learning rate,
optimizer, and augmentation---are as important as architectural choices for deep non-residual
CNNs, consistent with the broader findings of He et al.\ \cite{he2016deep} on deep network
optimization.

% ─────────────────────────────────────────────────────────────────────────────
\section{Method}
% ─────────────────────────────────────────────────────────────────────────────

\subsection{Dataset}

We use CIFAR-10 \cite{krizhevsky2009learning}, a standard benchmark consisting of 60,000
color images (32$\times$32 pixels) across 10 classes. The official training set (50,000
images) is split into 45,000 training samples and 5,000 validation samples using a fixed
seed (seed=42). The test set (10,000 images) is held out and evaluated only at the end of
each run. Images are normalized using the CIFAR-10 channel-wise mean $(0.4914, 0.4822, 0.4465)$
and standard deviation $(0.2023, 0.1994, 0.2010)$.

\subsection{Architectures}

\paragraph{ShallowCNN} consists of three convolutional blocks, each with a Conv2d (kernel
size 3, padding 1), BatchNorm2d, ReLU activation, and MaxPool2d(2) downsampling. Channel
widths are 32, 64, and 128. The classifier head is: Flatten $\to$ Linear(2048, 256) $\to$
ReLU $\to$ Dropout(0.5) $\to$ Linear(256, 10). Total parameters: \textbf{620,810}.

\paragraph{DeepCNN} follows a VGG-style design \cite{simonyanvery} with three blocks, each
containing two consecutive Conv-BN-ReLU layers followed by MaxPool2d(2) downsampling. Channel
widths are (64, 64), (128, 128), and (256, 256). The classifier head is: Flatten $\to$
Linear(4096, 512) $\to$ ReLU $\to$ Dropout(0.5) $\to$ Linear(512, 10). Total parameters:
\textbf{3,249,994}. Both models use He (Kaiming) initialization.

\subsection{Augmentation Levels}

Four augmentation pipelines are defined, applied only during training; evaluation always uses
the clean (no-augmentation) pipeline:

\begin{itemize}[noitemsep]
  \item \textbf{None:} Normalize only.
  \item \textbf{Basic:} RandomHorizontalFlip + Normalize.
  \item \textbf{Medium:} RandomHorizontalFlip + RandomCrop(32, padding=4) +
        ColorJitter(0.2, 0.2, 0.2, 0.1) + Normalize.
  \item \textbf{Aggressive:} Medium transforms + RandomGrayscale(p=0.1) +
        RandomErasing(p=0.5).
\end{itemize}

\subsection{Training Protocol}

All runs use cosine annealing LR schedule from $\mathrm{lr\_init}$ to 0 over 50 epochs;
batch size 128; weight\_decay=$10^{-4}$; seed 42. Three optimizer configurations were evaluated:

\begin{center}
\begin{tabular}{llllr}
\toprule
Run set & Optimizer & $\mathrm{lr\_init}$ & Architectures & \# Runs \\
\midrule
Original & SGD (mom=0.9) & 0.1 & ShallowCNN + DeepCNN & 8 \\
SGD rerun & SGD (mom=0.9) & 0.01 & DeepCNN only & 4 \\
Adam & Adam & $10^{-3}$ & DeepCNN / basic aug & 1 \\
\bottomrule
\end{tabular}
\end{center}

% ─────────────────────────────────────────────────────────────────────────────
\section{Experiments and Results}
% ─────────────────────────────────────────────────────────────────────────────

\subsection{Original 8 Runs (SGD, lr=0.1)}

\begin{table}[H]
\centering
\caption{Test accuracy for all 8 original runs (SGD, lr=0.1).}
\label{tab:original}
\begin{tabular}{llr}
\toprule
Architecture & Augmentation & Test Accuracy \\
\midrule
ShallowCNN & None       & 69.6\% \\
ShallowCNN & Basic      & \textbf{71.7\%} \\
ShallowCNN & Medium     & 68.8\% \\
ShallowCNN & Aggressive & 62.8\% \\
\midrule
DeepCNN    & None       & 41.5\% \\
DeepCNN    & Basic      & 10.4\% \; (\textit{collapsed}) \\
DeepCNN    & Medium     & 10.0\% \; (\textit{collapsed}) \\
DeepCNN    & Aggressive & 10.0\% \; (\textit{collapsed}) \\
\bottomrule
\end{tabular}
\end{table}

Figure~\ref{fig:learning_curves} shows the per-epoch learning curves for all 8 original runs.
ShallowCNN trained consistently across all augmentation levels (62.8--71.7\%). DeepCNN
showed a striking failure pattern: while it achieved 41.5\% without augmentation, all three
augmented conditions plateaued at approximately 10\%---equivalent to random chance for a
10-class problem---with training loss never decreasing.

\begin{figure}[H]
  \centering
  \includegraphics[width=\textwidth]{../cogs181_results/figures/fig1_learning_curves.png}
  \caption{Learning curves (train accuracy, dashed; val accuracy, solid) for all 8 original
  runs at SGD lr=0.1. The bottom row (DeepCNN) shows complete training collapse under any
  augmentation.}
  \label{fig:learning_curves}
\end{figure}

\begin{figure}[H]
  \centering
  \includegraphics[width=0.8\textwidth]{../cogs181_results/figures/fig2_bar_test_accuracy.png}
  \caption{Test accuracy for all original runs. ShallowCNN (blue) trains reliably; DeepCNN
  (red) collapses under augmentation at SGD lr=0.1.}
  \label{fig:bar_original}
\end{figure}

\begin{figure}[H]
  \centering
  \includegraphics[width=0.75\textwidth]{../cogs181_results/figures/fig3_augmentation_benefit.png}
  \caption{Augmentation benefit relative to the no-augmentation baseline, using the
  \emph{original} runs (SGD lr=0.1). ShallowCNN (blue) shows a mild $+$2.1\,pp gain from basic
  augmentation and a $-$6.7\,pp penalty from aggressive. DeepCNN (red) shows $-$31\,pp for all
  augmented conditions---a direct consequence of training collapse rather than augmentation
  sensitivity per se.}
  \label{fig:aug_benefit_orig}
\end{figure}

\subsection{LR Instability Diagnosis and Fix}

\textbf{Hypothesis:} $\mathrm{lr}=0.1$ is too large for a 6-layer non-residual network
combined with augmentation. Augmentation increases per-mini-batch input variance, which
amplifies gradient variance. In a shallow network this variance is tolerable; in a deep
non-residual network---where gradients must flow through many layers without skip
connections---the same variance tips training into an unstable regime.

\textbf{Fix:} Re-run all four DeepCNN conditions with $\mathrm{lr\_init}=0.01$.

\begin{table}[H]
\centering
\caption{DeepCNN test accuracy before and after LR reduction.}
\label{tab:lr_fix}
\begin{tabular}{lrrr}
\toprule
Augmentation & SGD lr=0.1 & SGD lr=0.01 & $\Delta$ \\
\midrule
None       & 41.5\% & 66.6\% & $+$25.1 pp \\
Basic      & 10.4\% & \textbf{74.0\%} & $+$63.6 pp \\
Medium     & 10.0\% & 65.1\% & $+$55.1 pp \\
Aggressive & 10.0\% & 44.3\% & $+$34.3 pp \\
\bottomrule
\end{tabular}
\end{table}

\begin{figure}[H]
  \centering
  \includegraphics[width=0.75\textwidth]{../cogs181_results_rerun/figures/fig_rerun_lr_comparison.png}
  \caption{DeepCNN test accuracy: original SGD lr=0.1 (dashed) vs.\ rerun SGD lr=0.01 (solid).
  Reducing the learning rate by 10$\times$ recovers DeepCNN from collapse.}
  \label{fig:lr_comparison}
\end{figure}

The hypothesis was confirmed. At lr=0.01, DeepCNN trained successfully across all augmentation
levels, with best accuracy of 74.0\% (basic augmentation).

\subsection{Augmentation Benefit: Properly Tuned Comparison}

\begin{table}[H]
\centering
\caption{Final test accuracies: ShallowCNN (SGD lr=0.1) vs.\ DeepCNN SGD rerun (lr=0.01).}
\label{tab:final}
\begin{tabular}{lrrrr}
\toprule
Architecture & None & Basic & Medium & Aggressive \\
\midrule
ShallowCNN (SGD lr=0.1)  & 69.6\% & \textbf{71.7\%} & 68.8\% & 62.8\% \\
DeepCNN (SGD lr=0.01)    & 66.6\% & \textbf{74.0\%} & 65.1\% & 44.3\% \\
\bottomrule
\end{tabular}
\end{table}

\begin{figure}[H]
  \centering
  \includegraphics[width=0.8\textwidth]{../cogs181_results_rerun/figures/fig_rerun_comparison.png}
  \caption{Side-by-side test accuracy: ShallowCNN (blue, SGD lr=0.1) vs.\ DeepCNN (red,
  SGD lr=0.01). Both peak at basic augmentation; DeepCNN drops more sharply under aggressive
  augmentation.}
  \label{fig:rerun_comparison}
\end{figure}

The key finding: basic augmentation provides a modest benefit for both architectures
(ShallowCNN: $+$2.1\,pp; DeepCNN: $+$7.4\,pp, both relative to each architecture's no-aug
baseline). As augmentation intensity increases, the architectures diverge sharply. DeepCNN
drops \textbf{29.7\,pp} from its best condition (basic: 74.0\%) to aggressive (44.3\%),
while ShallowCNN drops only \textbf{8.9\,pp} (71.7\% $\to$ 62.8\%), even after learning rate
correction. \textbf{ShallowCNN is substantially more robust to aggressive augmentation},
supporting the research hypothesis.

\subsection{Optimizer Comparison: Adam vs.\ SGD}

To assess whether SGD remained the limiting factor after the LR fix, we ran DeepCNN with
basic augmentation using Adam ($\mathrm{lr}=10^{-3}$, the standard default), holding all
other hyperparameters fixed.

\begin{table}[H]
\centering
\caption{Adam vs.\ SGD for DeepCNN / basic augmentation.}
\label{tab:optimizer}
\begin{tabular}{llr}
\toprule
Optimizer & $\mathrm{lr}$ & Test Accuracy \\
\midrule
SGD (momentum=0.9) & 0.01 & 74.0\% \\
\textbf{Adam} & $\mathbf{10^{-3}}$ & \textbf{89.1\%} \\
\bottomrule
\end{tabular}
\end{table}

\begin{figure}[H]
  \centering
  \includegraphics[width=0.5\textwidth]{../cogs181_results_rerun/figures/fig_adam_vs_sgd.png}
  \caption{Adam vs.\ SGD test accuracy for DeepCNN / basic augmentation. Adam outperforms
  the tuned SGD baseline by +15.1 percentage points.}
  \label{fig:adam_vs_sgd}
\end{figure}

Adam outperforms SGD by \textbf{+15.1 pp}---the single largest accuracy gain in the study.
This suggests that Adam's adaptive per-parameter learning rates substantially help in the
VGG-style deep non-residual setting, likely by compensating for the uneven gradient magnitude
distributions across deep layers that make a single fixed SGD learning rate suboptimal.

\subsection{Overfitting Analysis}

\begin{figure}[H]
  \centering
  \includegraphics[width=0.8\textwidth]{../cogs181_results/figures/fig4_train_val_gap.png}
  \caption{Train--validation accuracy gap (train$-$val at epoch 50) for all original runs.
  ShallowCNN gaps turn increasingly negative as augmentation intensifies (effective
  regularization). DeepCNN without augmentation shows a $-$2.7\,pp gap (val $>$ train,
  underfitting rather than overfitting). Augmented DeepCNN conditions show near-zero gaps
  ($+$0.6\,pp) because both train and val are stuck at $\sim$10\%.}
  \label{fig:train_val_gap}
\end{figure}

The train--validation accuracy gap (train$-$val at epoch 50) serves as a proxy for
overfitting. For ShallowCNN, the gap shrinks---and turns negative---as augmentation
intensifies: $+$3.8\,pp (none) $\to$ $-$0.1\,pp (basic) $\to$ $-$6.2\,pp (medium)
$\to$ $-$11.8\,pp (aggressive). Negative gaps are expected here: augmentation makes each
training image harder, suppressing training accuracy while the validation set remains clean.
This is evidence of effective regularization, not model failure. For original DeepCNN
(lr=0.1), the no-aug condition shows a $-$2.7\,pp gap (val $>$ train), indicating the model
is underfitting in training rather than overfitting. The three collapsed augmented conditions
show near-zero gaps ($+$0.6\,pp) because both train and val are stuck at random chance.

% ─────────────────────────────────────────────────────────────────────────────
\section{Discussion}
% ─────────────────────────────────────────────────────────────────────────────

\paragraph{Primary finding.} Basic augmentation provides consistent improvement for both
architectures. Aggressive augmentation hurts the deeper model disproportionately ($-30$ pp
vs.\ $-9$ pp). ShallowCNN is more robust to aggressive augmentation, supporting the hypothesis.

\paragraph{LR sensitivity.} DeepCNN fails completely at SGD lr=0.1 under augmentation.
Reducing to lr=0.01 recovers training. This is consistent with VGG-style networks requiring
more conservative SGD learning rates than shallow networks \cite{simonyanvery}.

\paragraph{Optimizer impact.} Switching from SGD (lr=0.01) to Adam (lr=$10^{-3}$) improves
DeepCNN by +15.1 pp on the best configuration, reaching 89.1\%. Adam's adaptive learning
rates appear especially beneficial for deep non-residual architectures, where layers develop
very different gradient magnitude distributions. Optimizer selection is at least as important
as learning rate tuning for this class of models.

\paragraph{Limitations.} The most significant limitation is the absence of residual
connections, which would likely change both LR sensitivity and augmentation robustness. The
Adam run was performed only for DeepCNN / basic augmentation; a full sweep across augmentation
levels would determine whether Adam also reduces the aggressive-augmentation penalty.
SGD and Adam learning rates were not jointly swept.

% ─────────────────────────────────────────────────────────────────────────────
\section{Conclusion}
% ─────────────────────────────────────────────────────────────────────────────

We conducted a systematic study of the interaction between CNN depth, augmentation intensity,
and optimizer choice on CIFAR-10. Our original 8-run experiment revealed an unexpected training
failure: DeepCNN collapses to random-chance accuracy under any augmentation at SGD lr=0.1.
Reducing the learning rate to 0.01 recovered training (best: 74.0\%, basic augmentation).
Switching to Adam ($\mathrm{lr}=10^{-3}$) further improved DeepCNN to \textbf{89.1\%}---the
highest accuracy in the study, and a +15.1 pp gain over the tuned SGD baseline.

Despite DeepCNN's higher peak accuracy when properly optimized, ShallowCNN is substantially
more robust to aggressive augmentation ($-9$ pp vs.\ $-30$ pp from best to aggressive
condition), supporting the research hypothesis. The overarching takeaways are: (1)
learning rate and augmentation must be tuned jointly for deep non-residual networks; (2)
Adam provides major gains over SGD for VGG-style architectures; (3) ShallowCNN provides
a more reliable baseline when training conditions cannot be carefully controlled. Future
work should apply Adam across all augmentation levels, explore residual connections, and
investigate weight decay as a complementary regularizer.

% ─────────────────────────────────────────────────────────────────────────────
\begin{thebibliography}{9}
% ─────────────────────────────────────────────────────────────────────────────

\bibitem{krizhevsky2009learning}
A.~Krizhevsky.
\newblock Learning multiple layers of features from tiny images.
\newblock Technical Report, University of Toronto, 2009.

\bibitem{simonyanvery}
K.~Simonyan and A.~Zisserman.
\newblock Very deep convolutional networks for large-scale image recognition.
\newblock In \textit{ICLR}, 2015.

\bibitem{he2016deep}
K.~He, X.~Zhang, S.~Ren, and J.~Sun.
\newblock Deep residual learning for image recognition.
\newblock In \textit{CVPR}, pp.~770--778, 2016.

\bibitem{shorten2019survey}
C.~Shorten and T.~M. Khoshgoftaar.
\newblock A survey on image data augmentation for deep learning.
\newblock \textit{Journal of Big Data}, 6(1):60, 2019.

\end{thebibliography}

\end{document}
